\renewcommand{\stagefigureprefix}{1}\setcounter{figure}{0}
\reporttitle{第一阶段：环境部署与向量执行过程}

本阶段对应赛题的环境部署验证、操作日志和向量加法执行过程，主要结果有三项：一是在WSL2中搭建了香山处理器仿真、NEMU参考模型和测试程序的完整构建与协同仿真环境；二是定位并修复了两端启动代码不一致导致的差分失败；三是三类vadd用例共24个元素在DiffTest下全部通过，并用波形观察了向量加法从执行到提交的过程。

实验分两次进行：10月9日完成环境搭建和首次验证，10月25日复核了英文hello输出、两种仿真模型的差分运行和含时钟的波形。第\ref{sec:1-index}、\ref{sec:1-logs}节集中列出版本信息和日志摘录，供复核使用。自定义vdot点积的设计与测试见第二阶段。

\section{协同验证框架与实验环境}

香山EMU、NEMU参考模型和裸机程序共同构成验证环境。EMU模拟硬件执行，NEMU按指令语义更新参考状态，DiffTest在指令提交时比较两者的架构状态。裸机程序提供可重复的输入和结果检查；共同启动flash负责在复位后跳转到程序入口。hello检查程序启动与串口输出，vadd则覆盖向量配置、加载、按元素加法、存储和提交。

硬件侧由Java与Mill解析Scala/Chisel工程，firtool处理硬件中间表示，Verilator与宿主C++编译器生成EMU。目标程序使用RISC-V GCC和binutils编译、链接并生成ELF、BIN及反汇编。NEMU在宿主侧构建为参考库，由运行参数显式加载。三条构建链的产物在同一次差分运行中汇合，构建记录和运行记录通过产物的SHA-256摘要对应起来。表\ref{tab:s1-artifacts}列出各类产物的作用，以及复核时需要查看的主要记录。

{\small\setlength{\tabcolsep}{3pt}\begin{longtable}{>{\raggedright\arraybackslash}p{0.19\linewidth}>{\raggedright\arraybackslash}p{0.34\linewidth}>{\raggedright\arraybackslash}p{0.39\linewidth}}
\caption{构建产物与复核入口}\label{tab:s1-artifacts}\\
\toprule
\textbf{产物} & \textbf{作用} & \textbf{主要记录}\\\midrule\endfirsthead
\caption[]{构建产物与复核入口（续）}\\
\toprule
\textbf{产物} & \textbf{作用} & \textbf{主要记录}\\\midrule\endhead
香山EMU & 执行DUT仿真 & 硬件源码、DUT配置、生成命令、编译工具与文件SHA-256\\
NEMU参考库 & 提供架构状态参考 & 固定源码、构建配置、实际加载路径与文件SHA-256\\
程序ELF/BIN & 执行hello和vadd用例 & ISA/ABI、入口、装载段、反汇编、镜像摘要\\
共同flash & 建立一致的启动输入 & 启动汇编、12字节镜像、反汇编与显式加载参数\\
运行记录 & 关联输入、结果及结束状态 & 原始日志、退出码、波形、事件CSV与信号索引\\
\bottomrule\end{longtable}}

图\ref{fig:S1-01}给出这一协同验证框架的整体结构。上半区为WSL2中的三条编译链：香山工程经JDK、Mill、firtool和Verilator生成EMU，AM程序与启动代码经RISC-V GCC生成共同的程序BIN和boot-flash，NEMU源码由宿主编译器构建为参考动态库。下半区为差分运行：EMU与NEMU加载同一镜像，DUT按架构提交推进参考模型，两端状态送入DiffTest比较，串口输出、波形、比较结果和退出状态统一归档。

\reportfigure{S1-01}{开发环境与软硬件协同验证框架}{10cm}{上半区展示WSL2中的三条编译链，下半区展示共同程序与启动输入驱动的差分运行。紫色虚线表示按架构提交推进参考模型；DUT与参考端的状态进入DiffTest，日志、波形和退出状态统一归档。}{2026-10-09历史环境的结构重绘；共同启动flash与程序BIN为分别记录的输入，固定版本见本章。}

10月9日实验在Windows下的WSL2 Ubuntu 22.04.5运行，宿主处理器为Ryzen 7 7840H，可用物理内存约27.7GiB。WSL配置20GB内存和16GB swap，构建与仿真放在WSL原生文件系统。DUT为MinimalConfig、单核、VLEN128，原波形构建采用\seqsplit{EMU\_THREADS=1}和\seqsplit{EMU\_TRACE=fst}。运行时显式指定128MB RAM、共同flash和固定参考库。

本阶段的源码版本、工具版本和10月25日复核的产物摘要集中列于本章末尾，两次实验所用的构建分开标注。

\section{部署与构建过程}

环境按本节步骤搭建。10月9日依次完成了NEMU参考库、香山EMU、测试程序和共同启动flash的构建，并完成首次差分运行，过程记录见表\ref{tab:s1-operations}。

\subsection{资源配置与工具位置}

WSL资源设置位于Windows用户目录的\seqsplit{\%UserProfile\%\textbackslash{}.wslconfig}：
\begin{Verbatim}[breaklines,breakanywhere,breaksymbolleft={},breaksymbolright={},fontsize=\footnotesize,baselinestretch=1]
[wsl2]
memory=20GB
swap=16GB
\end{Verbatim}

保存正在进行的WSL工作后，在PowerShell执行\Verb|wsl --shutdown|，重新进入Ubuntu并用\Verb|free -h|核对。原机器的交换文件位于\seqsplit{D:\textbackslash{}WSL\textbackslash{}wsl-swap.vhdx}，其他机器按本机磁盘布局配置。内存规划同时覆盖JVM、硬件生成工具、C++编译进程和文件缓存；频繁换页会增加构建时间，因此应根据可用物理内存调整JVM堆与编译并发。

Scala/Chisel、RISC-V程序、EMU/NEMU和波形处理都在WSL的Linux原生目录中执行。Windows目录下的工作树则用Windows原生Git管理：从WSL跨文件系统访问Windows目录很慢，实测一次Git检查约300秒仍未完成。源码与产物跨环境交换时，在接收端核对提交号和关键文件的SHA-256，运行输出使用独立目录。

\subsection{固定源码与工具安装}

在团队仓库执行\Verb|bash team/scripts/prepare-dependencies.sh|，按版本锁初始化xs-env、NEMU、nexus-am、DRAMsim3及递归依赖。团队仓库的香山子模块使用\Verb|git clone --recurse-submodules|恢复；NEMU的SimPoint按子模块记录初始化，其余下载依赖按锁定脚本检出。已有目录的版本不一致时，先保留现场并核对输入；每次构建前记录所用的分支、提交号和未提交改动，便于事后复现。

基础工具的安装入口如下。这里的命令用于恢复固定环境，安装后再核对实际工具版本。
\begin{Verbatim}[breaklines,breakanywhere,breaksymbolleft={},breaksymbolright={},fontsize=\footnotesize,baselinestretch=1]
source team/scripts/activate.sh
sudo bash "$TEAM_XS_ENV_ROOT/setup-tools.sh" --target base
sudo bash "$TEAM_XS_ENV_ROOT/setup-tools.sh" --target llvm
sudo bash "$TEAM_XS_ENV_ROOT/setup-tools.sh" --target mill
sudo apt-get install -y cmake ninja-build gtkwave
\end{Verbatim}

该入口来自固定xs-env的\href{https://github.com/OpenXiangShan/xs-env/blob/b0b506481e839a5df8a580e76f5aa246dc60fe7c/setup-tools.sh}{setup-tools.sh}。软件源中的补丁版本会变化，环境检查负责记录安装后实际使用的版本。Mill读取仓库的\seqsplit{.mill-version}选择项目版本，Scala、Chisel和firtool按项目解析结果核对。

xs-env原有的JDK安装入口指向Oracle的“21/latest”地址，版本会随时间变化。团队改为固定安装21.0.12 Linux x64归档，并用官方SHA-256校验，安装到用户目录：
\begin{Verbatim}[breaklines,breakanywhere,breaksymbolleft={},breaksymbolright={},fontsize=\footnotesize,baselinestretch=1]
export JAVA_HOME="$HOME/.local/share/graalvm-jdk-21.0.12"
bash team/scripts/install-graalvm.sh
source team/scripts/activate.sh
\end{Verbatim}

新终端沿用相同的JAVA\_HOME；脚本不会覆盖已有的其他版本。参考资料为\href{https://docs.oracle.com/en/graalvm/jdk/21/docs/getting-started/installation-linux/}{Oracle安装说明}、\href{https://download.oracle.com/graalvm/21/archive/graalvm-jdk-21.0.12\_linux-x64\_bin.tar.gz}{固定归档}和\href{https://download.oracle.com/graalvm/21/archive/graalvm-jdk-21.0.12\_linux-x64\_bin.tar.gz.sha256}{官方校验值}。

Verilator从固定提交构建，显式选择已安装的Clang，并按资源情况控制编译并发：
\begin{Verbatim}[breaklines,breakanywhere,breaksymbolleft={},breaksymbolright={},fontsize=\footnotesize,baselinestretch=1]
sudo apt-get install -y help2man perl autoconf flex bison libgoogle-perftools-dev \
  libjemalloc-dev numactl perl-doc zlib1g-dev liblz4-dev libfl2 libfl-dev
mkdir -p "$HOME/riscv-tools"
git clone https://github.com/verilator/verilator.git "$HOME/riscv-tools/verilator"
cd "$HOME/riscv-tools/verilator"
git checkout ea338be98e1e838d3518809ce8899f85a009963c
unset VERILATOR_ROOT
autoconf
./configure CC=/usr/local/llvm/bin/clang CXX=/usr/local/llvm/bin/clang++ LINK=/usr/local/llvm/bin/clang++
make -j2
sudo make install
\end{Verbatim}

安装后回到团队仓库根目录，运行\seqsplit{activate.sh}和\seqsplit{check-environment.sh}。原环境目录为\seqsplit{/home/riscv/riscv-lab/xs-env/XiangShan}；\seqsplit{activate.sh}将实际团队仓库设为NOOP\_HOME，依据TEAM\_XS\_ENV\_ROOT选择NEMU和AM，并设置\seqsplit{LINUX\_GNU\_TOOLCHAIN=1}。向量工作负载构建显式使用\Verb|MARCH=rv64gcv|。

\subsection{集成构建与操作记录}

确认源码、工具和输出目录后，依次构建参考库、EMU、目标程序及共同flash，再执行差分运行：
\begin{Verbatim}[breaklines,breakanywhere,breaksymbolleft={},breaksymbolright={},fontsize=\footnotesize,baselinestretch=1]
source team/scripts/activate.sh
bash team/scripts/check-environment.sh
bash team/scripts/stage1.sh nemu
bash team/scripts/stage1.sh emu
TEAM_ID=<团队或个人标识> bash team/scripts/stage1.sh workload
bash team/scripts/stage1.sh run
\end{Verbatim}

图\ref{fig:S1-02}展示上述命令对应的执行流程。环境准备、按序构建和DiffTest运行依次进行，验证通过后确认结果并归档证据；检查、构建或运行未通过时，分别定位配置与依赖、构建日志或运行差分的原因，修正后返回对应环节重新执行。每一步无论成功或失败，都保留命令、日志、退出码和SHA-256。

\reportfigure{S1-02}{环境部署与验证执行流程}{11cm}{依次完成环境准备、NEMU与EMU及程序构建、DiffTest运行和证据归档。检查、构建或运行失败时，按原因定位后返回对应环节；成功与失败记录均保留命令、退出码和摘要。}{2026-10-09部署记录与共同flash恢复过程的流程重绘。}

表\ref{tab:s1-operations}摘录10月9日构建后半段的主要操作及结果，时间为UTC+8。20:42时根据实际资源中断C++编译，把并行度调整为4后继续，已完成的目标文件全部保留。
{\small\setlength{\tabcolsep}{3pt}\begin{longtable}{>{\raggedright\arraybackslash}p{0.20\linewidth}>{\raggedright\arraybackslash}p{0.72\linewidth}}
\caption{2026年10月9日主要操作记录}\label{tab:s1-operations}\\
\toprule
\textbf{时间} & \textbf{操作与结果}\\\midrule\endfirsthead
\caption[]{2026年10月9日主要操作记录（续）}\\
\toprule
\textbf{时间} & \textbf{操作与结果}\\\midrule\endhead
20:42:15 & 中断并恢复C++构建，并行度调整为4，保留已完成对象。\\
20:53:32 & EMU构建完成，状态文件记录exit\_code=0。\\
20:55:26 & 工作负载与共同flash构建完成，记录\seqsplit{WORKLOAD\_BUILD\_PASS}。\\
21:00:28 & 正常差分运行结束，退出状态为0。\\
\bottomrule\end{longtable}}

这里的并行度只影响宿主机上的编译速度，与被测处理器的核数和EMU仿真线程数无关。

\section{共同启动输入与差分故障定位}

初次运行在PC\seqsplit{0x10000004}出现t0寄存器差异，记录的instrCnt为2，尚处于启动代码。比较默认flash和反汇编后发现，两端分别使用了3条与6条启动指令。修复由同一份\seqsplit{boot-flash.S}生成12字节\seqsplit{boot-flash.bin}，并通过运行参数显式提供给DUT与NEMU。随后两端以共同启动输入进入同一工作负载，在DiffTest启用的条件下完成hello/vadd检查。

\reportfigure{S1-07}{共同启动镜像与差分恢复}{12.1cm}{首次差异位于PC 0x10000004、instrCnt=2。共同启动文件由三条机器指令组成，图中0--3、4--7和8--11标示文件字节偏移。显式加载同一12字节flash后，三类共24个元素完成比较，并以GOOD TRAP、exit 0结束。}{2026-10-09原始失败和恢复日志、\seqsplit{boot-flash-disassembly.txt}、\seqsplit{hello-vadd-summary.log}；出处与摘要见\seqsplit{figures/evidence-20261105/figure-data.json}。}

图\ref{fig:S1-07}对照了修正前后的情况。上部表格逐字节列出共同启动文件：三条指令依次令t0=1、将其左移为0x80000000并跳转至该地址，0--3、4--7和8--11为文件内的字节偏移。下部左侧为修正前两端各用默认启动序列时的首次差异，右侧为两端显式加载同一文件后的结果：三类共24个元素通过比较，程序以GOOD TRAP和exit 0结束。

\section{hello输出与三类vadd结果}

\subsection{输入配置与结果解释}

vadd程序设置SEW8、LMUL1、VL8，执行两次\Verb|vle8.v|、一次\Verb|vadd.vv|和一次向量存储，再逐字节比较标量参考。三个用例检查基本加法、负数补码和8位回绕，合计比较24个活动元素。表\ref{tab:s1-vadd-values}给出三个用例的输出，分别以日志中的逐元素字节和波形中的低64位十六进制两种形式表示。

{\small\setlength{\tabcolsep}{3pt}\begin{longtable}{>{\raggedright\arraybackslash}p{0.18\linewidth}>{\raggedright\arraybackslash}p{0.43\linewidth}>{\raggedright\arraybackslash}p{0.31\linewidth}}
\caption{vadd输出的两种表示}\label{tab:s1-vadd-values}\\
\toprule
\textbf{用例} & \textbf{按元素顺序输出的字节} & \textbf{结果低64位}\\\midrule\endfirsthead
\caption[]{vadd输出的两种表示（续）}\\
\toprule
\textbf{用例} & \textbf{按元素顺序输出的字节} & \textbf{结果低64位}\\\midrule\endhead
正数相加 & 9 9 9 9 9 9 9 9 & \seqsplit{0x0909090909090909}\\
负数补码 & 254 252 250 248 246 244 242 240 & \seqsplit{0xf0f2f4f6f8fafcfe}\\
边界回绕 & 128 127 0 0 0 0 254 254 & \seqsplit{0xfefe000000007f80}\\
\bottomrule\end{longtable}}

每个元素按模256保留低8位。例如打印值254对应十六进制FE，按8位补码解释为$-2$。波形将多个字节合并显示为十六进制，右端为最低字节；日志按元素序号递增输出，因此两种书写顺序应按字节位置对应。VLEN128给出寄存器宽度，VL8给出本用例的活动元素数，表中结果覆盖低64位的八个元素。

\subsection{10月9日环境验证}

10月9日首次运行时团队标识尚未确定，hello程序使用临时标识\seqsplit{riscv-vdot-team}，因此输出中带有\texttt{identity\_status=placeholder}；显示的IPv4为\seqsplit{10.198.93.97}，对应当时WSL宿主。10月25日复核改用确定的标识重新构建，见下一小节。正常记录包含\seqsplit{HELLO\_VADD\_PASS cases=3 lanes=24}、DiffTest启用标记、GOOD TRAP和退出状态0；整程序统计为16557条指令、42962周期，覆盖启动、UART和结果校验。图\ref{fig:S1-03}为当次运行结束时的终端输出，可以看到身份行、三类vadd的PASS结果、DiffTest启用标记、GOOD TRAP以及退出状态0。

\reportfigure{S1-03}{hello与vadd正常结束输出}{8.5cm}{日志展示三类vadd结果、DiffTest运行和GOOD TRAP终态，输出中的身份与地址按原实验保留。}{用户截图\seqsplit{WindowsTerminal\_m5hXeDJnV1.png}及\seqsplit{evidence/hello-vadd-summary.log}；2026-10-09运行记录。}

\subsection{10月25日双模型复核}

10月25日（UTC+8）的复核以成员标识LiZufan001重新构建英文hello程序，输出格式为赛题要求的“hello xiangshan, I am xxx, IP address”，并标注地址来自构建宿主：
\begin{Verbatim}[breaklines,breakanywhere,breaksymbolleft={},breaksymbolright={},fontsize=\footnotesize,baselinestretch=1]
hello xiangshan, I am LiZufan001, IP address=172.27.184.54
identity_status=configured; IP address is build-host IPv4
\end{Verbatim}

同一BIN分别在封存生产EMU和独立FST模型上运行，均加载普通NEMU参考库并启用DiffTest。两次运行都通过三组、24个元素检查，命中GOOD TRAP并正常退出；每次为16621条指令、43249周期。这里的周期同样对应完整hello/vadd程序。表\ref{tab:s1-t04-runs}列出两次运行的实际起止时间与结果，时间取自各自的运行回执。

{\small\setlength{\tabcolsep}{3pt}\begin{longtable}{>{\raggedright\arraybackslash}p{0.22\linewidth}>{\raggedright\arraybackslash}p{0.38\linewidth}>{\raggedright\arraybackslash}p{0.32\linewidth}}
\caption{10月25日双模型运行记录（UTC+8）}\label{tab:s1-t04-runs}\\
\toprule
\textbf{模型} & \textbf{实际起止时间} & \textbf{运行结果}\\\midrule\endfirsthead
\caption[]{10月25日双模型运行记录（UTC+8）（续）}\\
\toprule
\textbf{模型} & \textbf{实际起止时间} & \textbf{运行结果}\\\midrule\endhead
生产EMU & 01:09:21--01:10:18 & DiffTest通过，GOOD TRAP，exit 0\\
独立FST模型 & 01:28:57--01:31:24 & DiffTest通过，GOOD TRAP，exit 0\\
\bottomrule\end{longtable}}

FST波形模型复用已有构建的1807个生成文件单独编译；用发布的源码重新编译得到的程序BIN和flash，与实际运行时的输入逐字节相同。对应的\href{https://github.com/cc77-builder/riscv/blob/2690c7370fb7c94111b4619bbac352e9be4ccc4d/team/evidence/stage1-current-20261025/acceptance.json}{验收记录}、\href{https://github.com/cc77-builder/riscv/blob/2690c7370fb7c94111b4619bbac352e9be4ccc4d/team/evidence/stage1-current-20261025/production-run/run.json}{生产运行回执}和\href{https://github.com/cc77-builder/riscv/blob/2690c7370fb7c94111b4619bbac352e9be4ccc4d/team/evidence/stage1-current-20261025/trace-run/run.json}{FST运行回执}见团队仓库。

\section{向量执行与波形观察}

向量指令经取指、译码、重命名与调度进入执行单元。两个\Verb|vle8.v|提供输入，\Verb|vadd.vv|产生结果，随后经过写回与ROB提交。本文用提交PC\seqsplit{0x80000138}、机器字\seqsplit{0x021101d7}及\Verb|vecwen=1, wdest=3|关联三类用例，并分别记录结果输出和架构提交两个观察点。

\subsection{10月9日事件与原始截图}

表\ref{tab:s1-original-events}沿用原FST提取CSV的事件坐标，列出三类用例在执行输入、结果输出和ROB提交三处的位置以及输出低64位。局部截图未显示时钟；截图中的ps只是波形工具的显示单位，数值为仿真记录的时间坐标，不代表实际时间。本表只用于定位事件，指令时延与性能在第三阶段用专门的测量窗口评估。

{\small\setlength{\tabcolsep}{3pt}\begin{longtable}{>{\raggedright\arraybackslash}p{0.15\linewidth}>{\raggedright\arraybackslash}p{0.16\linewidth}>{\raggedright\arraybackslash}p{0.16\linewidth}>{\raggedright\arraybackslash}p{0.16\linewidth}>{\raggedright\arraybackslash}p{0.29\linewidth}}
\caption{10月9日vadd事件坐标}\label{tab:s1-original-events}\\
\toprule
\textbf{用例} & \textbf{输入} & \textbf{输出} & \textbf{提交} & \textbf{输出低64位}\\\midrule\endfirsthead
\caption[]{10月9日vadd事件坐标（续）}\\
\toprule
\textbf{用例} & \textbf{输入} & \textbf{输出} & \textbf{提交} & \textbf{输出低64位}\\\midrule\endhead
正数 & 19625 & 19626 & 19631 & \seqsplit{0x0909090909090909}\\
负数 & 25460 & 25461 & 25466 & \seqsplit{0xf0f2f4f6f8fafcfe}\\
边界 & 33009 & 33010 & 33015 & \seqsplit{0xfefe000000007f80}\\
\bottomrule\end{longtable}}

图\ref{fig:S1-06}把表\ref{tab:s1-original-events}的坐标画成时间线。三类用例都在输入后1个周期给出结果，再经5个周期完成ROB提交；每个用例的三个事件都属于同一条vadd指令。

\reportfigure{S1-06}{vadd输入、输出与提交事件时间线}{8.2cm}{三类用例按指令身份关联执行端输入、结果输出和ROB提交；数值沿用事件CSV坐标。}{\seqsplit{evidence/vadd-commit-events.csv}与\seqsplit{evidence/vector-fix-alu-events.csv}；2026-10-09事件重绘，生成数据见\seqsplit{figures/evidence-20261105}。}

图\ref{fig:S1-04}是正数用例的原始GTKWave截图：子图（a）的游标位于19626处的输出事件，三个事件的位置同时可见；子图（b）、（c）分别给出19625输入游标和19631提交游标处的信号值。图\ref{fig:S1-05}为负数与边界用例的输出波形，两者的输出低64位分别为\seqsplit{F0F2F4F6F8FAFCFE}和\seqsplit{FEFE000000007F80}，与表\ref{tab:s1-vadd-values}一致。

\reportfigure{S1-04}{vadd正数用例的执行与提交波形}{16.0cm}{(a)输出事件及相关位置；(b)输入游标处信号值；(c)提交游标处信号值。19625、19626和19631对应表中的三个事件坐标。}{用户上传的正数输出截图及\seqsplit{riscv-lab/picture}中的输入、提交截图；2026-10-09原基线波形，保留原始图并裁剪外部空白。}

\reportfigure{S1-05}{vadd负数与边界用例结果波形}{10.0cm}{负数和边界样例输出低64位分别为\seqsplit{F0F2F4F6F8FAFCFE}与\seqsplit{FEFE000000007F80}，游标对应各自输入、输出和提交位置。}{用户上传\seqsplit{msrdc\_yWBpPeYM8q.png}与\seqsplit{msrdc\_6cIzJ9aIVJ.png}；同一2026-10-09波形的八个活动元素。}

原完整FST路径记录为\seqsplit{/home/riscv/riscv-lab/stage1-evidence/hello-vadd.fst}，大小220295598字节。复核可先从提交PC定位同次反汇编，再根据CSV打开波形窗口，最后联合查看时钟、valid/ready、操作数、结果及事务标识。事件CSV和波形提取记录见团队仓库的\seqsplit{team/evidence/stage1-20261009/}。

\subsection{10月25日含时钟的调试记录}

10月25日的FST波形模型保留了实际时钟信号。从中选取两个向量ALU（VIAluFix）实例和八个提交槽，共127个信号，逐个检查33249个时钟上升沿，得到九条事件：三次FU输入、三次FU输出和三次DPI-C提交观测。三笔vadd都在\seqsplit{exus\_2}中的VIAluFix执行。

采样取每个上升沿之前、时钟低电平期间的稳定值，周期字段取当时的\seqsplit{DifftestTrapEvent.io\_cycleCnt}，FST原始时戳另存。DPI-C是仿真模型调用C函数的接口，提交记录经过\seqsplit{Rob.scala}中\Verb|delay=3|的观察路径，表内DPI-C列因此表示提交接口的观测时刻。按上述规则得到的九条事件列于表\ref{tab:s1-t04-events}，表中同时给出每笔运算的ROB flag、value和目的物理寄存器。

{\small\setlength{\tabcolsep}{3pt}\begin{longtable}{>{\raggedright\arraybackslash}p{0.10\linewidth}>{\raggedright\arraybackslash}p{0.16\linewidth}>{\raggedright\arraybackslash}p{0.16\linewidth}>{\raggedright\arraybackslash}p{0.17\linewidth}>{\raggedright\arraybackslash}p{0.33\linewidth}}
\caption{10月25日按时钟沿采样的事件记录}\label{tab:s1-t04-events}\\
\toprule
\textbf{用例} & \textbf{FU输入timer} & \textbf{FU输出timer} & \textbf{DPI-C观测timer} & \textbf{ROB flag/value/pdest}\\\midrule\endfirsthead
\caption[]{10月25日按时钟沿采样的事件记录（续）}\\
\toprule
\textbf{用例} & \textbf{FU输入timer} & \textbf{FU输出timer} & \textbf{DPI-C观测timer} & \textbf{ROB flag/value/pdest}\\\midrule\endhead
正数 & 19525 & 19526 & 19531 & 1 / 13 / 51\\
负数 & 25526 & 25527 & 25532 & 0 / 21 / 57\\
边界 & 33193 & 33194 & 33199 & 1 / 25 / 63\\
\bottomrule\end{longtable}}

三类数值结果与表\ref{tab:s1-vadd-values}一致。表中的timer是周期计数。10月25日的程序换用新标识重新编译，指令地址和执行时刻都与10月9日不同，因此表\ref{tab:s1-t04-events}与表\ref{tab:s1-original-events}的绝对坐标不同；但每个用例从输入到输出相隔1个周期、从输出到提交相隔5个周期，两次完全一致。完整采样定义、信号清单、原始时戳及摘要保存在\href{https://github.com/cc77-builder/riscv/blob/2690c7370fb7c94111b4619bbac352e9be4ccc4d/team/evidence/stage1-current-20261025/wave/wave.json}{wave.json}；只含所选信号的FST可直接用GTKWave查看。接口与观察路径分别对应固定源码的\href{https://github.com/cc77-builder/riscv/blob/34bae6d4acceac939f958199dae53c14bc1f6183/src/main/scala/xiangshan/backend/fu/wrapper/VIAluFix.scala}{VIAluFix包装层}和\href{https://github.com/cc77-builder/riscv/blob/34bae6d4acceac939f958199dae53c14bc1f6183/src/main/scala/xiangshan/backend/rob/Rob.scala}{Rob.scala}。

提取工具本身也做了24项检查，覆盖日志完整性、退出状态、输入文件摘要和关键结果，并用缺少FU事件、缺半个周期、操作数未知和事务标识错误等反例确认工具能发现问题。

\section{版本与复现索引}\label{sec:1-index}

本节列出两次实验所用的源码、工具和产物，供复现和核对使用。

\subsection{源码与工具版本}

表\ref{tab:s1-source-lock}列出10月9日部署时固定的源码版本，完整递归依赖见团队仓库的\seqsplit{team/evidence/stage1-20261009/submodules-after-build.lock.txt}。
{\small\setlength{\tabcolsep}{3pt}\begin{longtable}{>{\raggedright\arraybackslash}p{0.23\linewidth}>{\raggedright\arraybackslash}p{0.69\linewidth}}
\caption{10月9日源码版本锁}\label{tab:s1-source-lock}\\
\toprule
\textbf{组件} & \textbf{提交号}\\\midrule\endfirsthead
\caption[]{10月9日源码版本锁（续）}\\
\toprule
\textbf{组件} & \textbf{提交号}\\\midrule\endhead
xs-env & \seqsplit{b0b506481e839a5df8a580e76f5aa246dc60fe7c}\\
香山基线 & \seqsplit{16ae9ddcda54fc9a2fddffad73174cf793ac7814}\\
NEMU & \seqsplit{a48a11f208ff81ad765591d83271d1b339fb3861}\\
nexus-am & \seqsplit{086a238062fdfb3281b745a537b361983d0154b5}\\
DRAMsim3 & \seqsplit{1d2a9bf6da8e050d975ef2669c42faa93943f489}\\
DiffTest & \seqsplit{62e4540c03c03aa947bdb67f4b6bf66f1a29c5d0}\\
yunsuan & \seqsplit{cadd3c2f43096e253d61296b33ac697be8354e29}\\
\bottomrule\end{longtable}}

{\small\setlength{\tabcolsep}{3pt}\begin{longtable}{>{\raggedright\arraybackslash}p{0.34\linewidth}>{\raggedright\arraybackslash}p{0.58\linewidth}}
\caption{10月9日工具版本}\label{tab:s1-tools}\\
\toprule
\textbf{工具} & \textbf{版本}\\\midrule\endfirsthead
\caption[]{10月9日工具版本（续）}\\
\toprule
\textbf{工具} & \textbf{版本}\\\midrule\endhead
Oracle GraalVM JDK / javac & 21.0.12\\
LLVM / Clang & 19.1.7\\
Mill & 项目版本0.12.3\\
Scala / Chisel & 2.13.15 / 6.6.0\\
firtool & 1.62.1\\
Verilator & 5.052，提交\seqsplit{ea338be98e1e838d3518809ce8899f85a009963c}\\
RISC-V GCC / GNU binutils & 11.4.0（\seqsplit{riscv64-linux-gnu-}前缀）/ 2.38\\
宿主GCC / G++ & 11.4.0\\
Make / CMake / Ninja & 4.3 / 3.22.1 / 1.10.1\\
Python / GTKWave & 3.10.12 / 3.3.104\\
\bottomrule\end{longtable}}

表\ref{tab:s1-tools}为同一环境实际使用的工具版本，完整版本输出和Ubuntu软件包清单见同一目录的\seqsplit{tool-versions.txt}与\seqsplit{ubuntu-packages.lock.tsv}。

\subsection{10月25日复核的产物摘要}

10月25日复核所用硬件构建的源码提交为\seqsplit{34bae6d4acceac939f958199dae53c14bc1f6183}。表\ref{tab:s1-t04-identities}列出这次运行主要产物的SHA-256，与团队仓库\seqsplit{team/evidence/stage1-current-20261025/}中的运行和验收记录一致。
{\footnotesize\setlength{\tabcolsep}{3pt}\begin{longtable}{>{\raggedright\arraybackslash}p{0.23\linewidth}>{\raggedright\arraybackslash}p{0.69\linewidth}}
\caption{10月25日复核主要产物的SHA-256}\label{tab:s1-t04-identities}\\
\toprule
\textbf{对象} & \textbf{SHA-256}\\\midrule\endfirsthead
\caption[]{10月25日复核主要产物的SHA-256（续）}\\
\toprule
\textbf{对象} & \textbf{SHA-256}\\\midrule\endhead
生产EMU & \seqsplit{ff840528eb1009e94906eefb3773ca321050f309553543ab26f0db23c410e30f}\\
FST模型 & \seqsplit{1ff3f457ff582d0c3155a15dafa1fa60b8275d9356e29a0b7806c9d3915ce9f7}\\
普通NEMU参考库 & \seqsplit{8af118e5519db7ef6001010a2d547707c5d47630412e80abfc3f343bfef987c9}\\
共同程序BIN & \seqsplit{d81437b44473cef2b0c5777330e13fd076567d4fc1a4296be35d971e51a4dcad}\\
共同flash & \seqsplit{a3360212cf7bb784119c1719c9258362f8567803601b2359aee2c104864fa994}\\
完整FST（181143183字节） & \seqsplit{cf8abe9104fb955c7dc4bf1efda6588a61060fce33857b9ddc6d5020421be299}\\
只含所选信号的FST（393288字节） & \seqsplit{4c9f442704871010026654fd9d60854024b270a83be6354be5dee3ff25df24e1}\\
\bottomrule\end{longtable}}

\subsection{文件与记录索引}

运行源码位于团队仓库的\seqsplit{team/workloads/hello-vadd}。10月9日的证据在\seqsplit{team/evidence/stage1-20261009}，10月25日复核的证据在\seqsplit{team/evidence/stage1-current-20261025}；第\ref{sec:1-logs}节摘录的日志另存于报告工程的\seqsplit{evidence/}目录。完整FST文件较大，按记录中的路径和SHA-256定位。各记录的用途见表\ref{tab:s1-evidence-index}。

{\small\setlength{\tabcolsep}{3pt}\begin{longtable}{>{\raggedright\arraybackslash}p{0.50\linewidth}>{\raggedright\arraybackslash}p{0.42\linewidth}}
\caption{第一阶段记录索引}\label{tab:s1-evidence-index}\\
\toprule
\textbf{记录} & \textbf{用途}\\\midrule\endfirsthead
\toprule
\textbf{记录} & \textbf{用途}\\\midrule\endhead
\multicolumn{2}{l}{报告工程的\texttt{evidence/}目录：}\\
\seqsplit{workload-build.log} & 工作负载与flash的构建命令及成功状态（第\ref{sec:1-logs}节摘录）\\
\seqsplit{hello-vadd-summary.log} & 10月9日运行摘要（第\ref{sec:1-logs}节摘录）\\
\seqsplit{rendered/hello-vadd-production-20261025.txt}与\seqsplit{hello-vadd-trace-20261025.txt} & 10月25日两次运行的终端输出（第\ref{sec:1-logs}节摘录）\\
\multicolumn{2}{l}{团队仓库的\texttt{team/evidence/stage1-20261009/}：}\\
\seqsplit{tool-versions.txt}、\seqsplit{ubuntu-packages.lock.tsv} & 工具版本原始输出与软件包清单\\
\seqsplit{submodules-after-build.lock.txt} & 递归依赖完整提交锁\\
\seqsplit{vadd-commit-events.csv}、\seqsplit{vector-fix-alu-events.csv} & 提交记录与三组向量ALU输入输出记录，表\ref{tab:s1-original-events}的来源\\
\seqsplit{wave-extraction.log} & 10月9日的波形事件提取记录\\
\multicolumn{2}{l}{团队仓库的\texttt{team/evidence/stage1-current-20261025/}：}\\
\seqsplit{production-run/run.json}、\seqsplit{trace-run/run.json} & 10月25日两次运行的命令、输入摘要、退出状态和日志摘要\\
\seqsplit{wave/wave.json} & 采样定义、127个信号和九条事件\\
\seqsplit{acceptance.json}、\seqsplit{RESULT.md} & 输入等价、24项工具检查与复现说明\\
\bottomrule\end{longtable}}

复核时先检查运行回执和通过摘要，再核对程序、参考库与flash的SHA-256，最后用反汇编和事件CSV定位波形。详细操作入口见\href{https://github.com/cc77-builder/riscv/blob/eca41d7cd959f5f3c4f240e382919c492d1ccb96/team/docs/stage1.md}{阶段一记录}、\href{https://github.com/cc77-builder/riscv/blob/347b1ae7eafb0f30c3b2cc1f8da8d9c8cdd3f98e/team/docs/environment.md}{环境部署方法}和\href{https://github.com/cc77-builder/riscv/blob/2690c7370fb7c94111b4619bbac352e9be4ccc4d/team/evidence/stage1-current-20261025/RESULT.md}{10月25日复核的复现说明}。

\section{操作日志摘录}\label{sec:1-logs}

本节摘录主要的构建和运行输出；版本锁、事件CSV和运行回执见表\ref{tab:s1-evidence-index}所列位置。

\subsection{10月9日工作负载构建}
\VerbatimInput[firstline=5,lastline=17,breaklines,breakanywhere,breaksymbolleft={},breaksymbolright={},fontsize=\footnotesize,baselinestretch=1]{evidence/workload-build.log}
\VerbatimInput[firstline=22,lastline=26,breaklines,breakanywhere,breaksymbolleft={},breaksymbolright={},fontsize=\footnotesize,baselinestretch=1]{evidence/workload-build.log}

\subsection{10月9日正常运行摘要}
\VerbatimInput[breaklines,breakanywhere,breaksymbolleft={},breaksymbolright={},fontsize=\footnotesize,baselinestretch=1]{evidence/hello-vadd-summary.log}

\subsection{10月25日生产模型输出}
\VerbatimInput[breaklines,breakanywhere,breaksymbolleft={},breaksymbolright={},fontsize=\footnotesize,baselinestretch=1]{evidence/rendered/hello-vadd-production-20261025.txt}

\subsection{10月25日FST模型输出}
\VerbatimInput[firstline=11,lastline=21,breaklines,breakanywhere,breaksymbolleft={},breaksymbolright={},fontsize=\footnotesize,baselinestretch=1]{evidence/rendered/hello-vadd-trace-20261025.txt}
